\chapter*{Preface}
In $\Z[\sqrt{-5}]$ the number $6$ is both $2\cdot3$ and $(1+\sqrt{-5})(1-\sqrt{-5})$, and unique factorisation is lost.
Kummer and Dedekind found it again by factoring \emph{ideals} instead of numbers, and with that repair a whole world
opened. This book is a journey through that world: the arithmetic of number fields, of the $p$-adic numbers, and of
curves, up to the modular forms that tie them together.

The first part studies number fields. Their integers form a lattice whose volume, the discriminant, decides which primes
ramify. Every prime number splits into prime ideals in a way that a single polynomial modulo $p$ reveals, and in Galois
extensions this splitting is governed by one group element, the Frobenius. Minkowski's convex body theorem shows that
class groups are finite and that unit groups are as large as they can be, and Kummer's cyclotomic fields give the
first serious attack on Fermat's Last Theorem.

The second part looks at one prime at a time and then at all of them together. Extensions of the $p$-adic numbers are
classified by two numbers, $e$ and $f$. The local--global principle of Hasse and Minkowski decides every quadratic
equation by congruences and inequalities, and settles which integers are sums of three squares. Zeta functions count
ideals, compute class numbers, and prove Chebotarev's theorem that every conjugacy class of a Galois group occurs as a
Frobenius with the frequency its size predicts. Ad\`eles put the real and $p$-adic worlds on an equal footing, and
Tate's thesis derives the functional equation of the zeta function from Fourier analysis on them. Class field theory
then describes all abelian extensions of $\Q$: the Kronecker--Weber theorem is proved in full, and the two main theorems
of the general theory, Artin reciprocity and the existence theorem, are stated, clearly marked as admitted, and used to
decide which primes are of the form $x^2+ny^2$.

The third part is arithmetic geometry. Plane curves meet as Bézout predicts, and Riemann--Roch measures their functions.
On an elliptic curve the points form a group, finitely generated over $\Q$ by the Mordell--Weil theorem, and over a
finite field their number obeys Hasse's bound, a special case of the Riemann hypothesis for curves that we prove with
Bombieri and Stepanov's elementary method. The book ends with modular forms: the discriminant function
$\Delta=q\prod(1-q^n)^{24}$, its coefficients $\tau(n)$, Hecke operators and the multiplicativity they explain,
Hecke's proof that their $L$-functions have a functional equation, and a form whose coefficients count the points of an
elliptic curve modulo every prime.

The aim throughout is to understand why the theorems are true, and to see them at work. Every result is proved, apart from the two admitted theorems of class field theory. Pause
and think boxes ask you to guess before reading on. Each chapter ends with exercises, graded from warm-ups
($\star$) to small gems ($\star\star\star$), each with a hint and a complete solution that first explains how to find
the idea and then writes it up. A companion Python notebook for each chapter computes with PARI/GP, checks the numbers
quoted in the text, and opens a window on computer science: factoring polynomials over finite fields, lattice
reduction and lattice cryptography, number-theoretic transforms, elliptic-curve cryptography, point counting, codes from
curves, and the series that compute $\pi$ to millions of digits.

The beauty of the subject is in its coincidences that turn out not to be coincidences. Whether a prime $p\equiv1\pmod 3$ splits completely in
$\Q(\sqrt[3]2)$ is decided by writing $p=x^2+27y^2$. The number of points of $y^2=x^3-x$ modulo $p$ is predicted by
Gaussian integers. The coefficients of a power series in $q$ count the points of a curve. Each of these is a shadow of
one large structure, and the chapters that follow try to make that structure visible.
